% !TeX spellcheck = <none>
\documentclass[12pt]{scrartcl}
\usepackage[english]{babel}
\usepackage[utf8]{inputenc}
\usepackage{csquotes}
\usepackage[T1]{fontenc}
\usepackage{lmodern}
\usepackage{geometry}
\newgeometry{left=3cm, right=3cm, top=3cm, bottom=3cm}
\usepackage[onehalfspacing]{setspace}
\usepackage{graphicx}
\usepackage{caption}
\usepackage{subcaption}

\usepackage{textcomp}         % degree symbol etc.
\usepackage{newunicodechar}   % map any stray Unicode code points

\usepackage{amsmath}   % gives \tfrac, \dfrac, \eqref, etc.
\usepackage{booktabs}   % for \toprule, \midrule, \bottomrule
\usepackage{enumitem}   % enables [leftmargin=...] in itemize/description



\usepackage[backend=biber, style=authoryear, sorting=nty]{biblatex}

\addbibresource{Bib.bib}

\usepackage{mathptmx}
\usepackage{times}

\usepackage{multicol}
\setlength{\columnsep}{1cm}

\usepackage{fancyhdr}
\pagestyle{fancy}
\fancyhf{}
\lhead{}
\rhead{Levin Damian Antonin Giersch}
\rfoot{Page \thepage}



\begin{document}
	\begin{titlepage}
		\centering
		
		\noindent
		\textbf{\huge A harmonised point-extraction tool for heterogeneous paleoclimate NetCDF datasets}\\
		\vspace{2cm}
		\textbf{\large Essay or reort or something for geoinformatics project seminar}\\
		\vspace{8cm}
		\begin{tabular}{ll}
			\textbf{Name:} & Levin Giersch \\
			\textbf{Student ID:} & 1690 0515040 \\
			\textbf{E-Mail:} & l.giersch@campus.tu-berlin.de\\
			\textbf{Faculty:} & Faculty VI - Planning Building Environment\\
			\textbf{Program:} & Geodesy and Geoinformation Science, M.Sc. \\
			\textbf{Date:} & \today \\
			\textbf{Lecturer:} & Prof. Dr.-Ing. Martin Kada
		\end{tabular}
		
		\newpage
		\tableofcontents
		\thispagestyle{empty}
		
	\end{titlepage}
	
	
	
	%%%%%%%%%%%%%%%%%%%%%%%%%%%%%%%%%%%%%%%%%%%%%%%%%%%%%%%%%%%%%%%%%%%%%%%%%%%%%%%%%%%%%%%%%%%%%%%%%%%%
	%\newpage
	\setcounter{page}{1}
	\section*{Abstract}
	\addcontentsline{toc}{section}{Abstract}
	
	\begin{abstract}
	
	\end{abstract}
	
	%-------------------------------------------------------------

	\section{Intro}
	- Paleoclimate models are useful for:\\
		- Reconstructing Earth's climate before direct measurements existed\\
		- Understanding how the climate responds to forcings like solar output, volcanic eruptions, orbital cycles, and CO2\\
		- Testing and validating modern climate models against known past states\\
		- Establishing the natural variability range of climate, separate from human influence\\
		- Constraining climate sensitivity (temperature change per CO2 doubling)\\
		- Studying planetary climate evolution — Mars, Venus, and other terrestrial bodies\\
		- Using past warm periods (e.g. Pliocene) as analogues for near-future climate projections\\
	
	- there are many different paleoclimate models.\\
	- coparing them is important to detect outliers, irregularities or strong differences eg. to help improve the models\\
	
	
	\section{The Data}
	- can be downloaded via data\_downloader.ipynb
	
	\subsection{Trace21k}
	- Simulation of the Transient Climate of the Last 21,000 Years\\
	- Chosen Version: TraCE Main Land Decadal Average\\
	- developed by: see source \\
	- year of development: see source\\
	- how where they made: fully-coupled, non-accelerated atmosphere-ocean-sea ice-land surface CCSM3 simulation\\
	- what was the goal: \\
	- nr. files: 2\\
	- size (gb): 449,2 MB\\
	- type: netcdf\\
	- used variables: TSA, TSOI\\
	- how to get: downloaded via wget python script\\
	- other info:  \\
	- extend time: 22,000 BCE to 1990 CE \\
	- resolution time: decadal averages\\
	- extend space: global land\\
	- resolution space: T31\_gx3\\
	- source: \cite{trace21k_data}\\

	\subsection{CHELSA-TraCE21k}
	- Chosen Version: CHELSA-TraCE21k-centennial\\
	- developed by: see source\\
	- year of development: see source \\
	- how where they made: generated with the CHELSA-TraCE21k downscaling model. \\
	- what was the goal: \\
	- nr. files: \\
	- size (gb): \\
	- type: geotiff -> need to be converted to netcdf\\
	- used variables: tasmin, tasmax -> tasmin+tasmax/2 to get mean\\
	- how to get: downloaded via wget python script \\
	- other info: \\
	- extend time: Start Time 21k BP, End Time 0 BP\\
	- resolution time: monthly \\
	- extend space: global \\
	- resolution space: kilometer scale\\
	- source: \cite{chelsa_trace21k_paper}, \cite{chelsa_trace21k_data} \\
	
	
	\subsection{Beyer 2021 (v4)}
	- Chosen Version: 12293345 version 4 (figshare, published 24 June 2021)\\
	- developed by: Robert M. Beyer, Mario Krapp \& Andrea Manica (Department of Zoology, University of Cambridge)\\
	- year of development: original paper published 14 July 2020 (received 03 July 2019); dataset v4 published 24 June 2021; addendum paper published 29 September 2021\\
	- how were they made: Delta-method-based downscaling and bias correction, combining medium-resolution HadCM3 transient simulations with high-resolution HadAM3H snapshots and present-day observational data\\
	- what was the goal: Fill the gap between high-temporal-resolution but low-spatial-resolution palaeoclimate simulations and high-spatial-resolution snapshot-only datasets\\
	- nr. files: main NetCDF data file plus R, Python and MATLAB example scripts (verify exact count on the figshare page)\\
	- size (GB): 1.24 GB (main NetCDF data file)\\
	- type: NetCDF (data) + R / Python / MATLAB scripts\\
	- used variables: \\
	- how to get: downloaded via wget Python script\\
	- other info: v4 adds global monthly NPP relative to v3; documented in the 2021 addendum\\
	- extend time: 120{,}000 BP to pre-industrial modern era (0 BP)\\
	- resolution time: 2000-yr steps from 120000--22000 BP; 1000-yr steps from 22000 BP to pre-industrial (72 snapshots total)\\
	- extend space: global terrestrial (full 360° longitude; latitude coverage 720 $\times$ 300 cells at 0.5°, i.e. 150° of latitude)\\
	- resolution space: 0.5° equirectangular grid (\textasciitilde55 km at the equator)\\
	- source: \cite{beyer2020_paper}, \cite{beyer2021_addendum}, \cite{beyer2021_data_4}\\
	
	
	\subsection{PalMod 2}
	- Chosen Version: \\
	- developed by: \\
	- year of development: \\
	- how where they made: \\
	- what was the goal: \\
	- nr. files: \\
	- size (gb): \\
	- type: \\
	- used variables: \\
	- how to get: \\
	- other info: \\
	- extend time: \\
	- resolution time: \\
	- extend space: \\
	- resolution space: \\
	- source: 	\\
	
	
	%-------------------------------------------------------------
	\section{Homogenization}
	\subsection{Variables}
	- which do i choose, are they directly comparable? \\
	
	\subsection{Time}
	- time extend \\
	- time resolution \\
	
	\subsection{Space}
	- grid sizes \\
	- changes in north south \\
	- different resolutions for different times \\
	
	%-------------------------------------------------------------
	\section{Point Extraction}
	- do i just take the grid cell the point is in? \\
	
	%-------------------------------------------------------------
	\section{Results)}
	- For one example Point \\
	- differences between models \\
	
	%-------------------------------------------------------------
	\section{Future}
	- Auto iterate over points and compare \\
	- correlate with eg. north south, climate zones, water vs land etc. \\
	
	
	%-------------------------------------------------------------
	\section{Conclusion}
	- hard to compare due to differences eg. resolutions \\
	- comparison can helt to find out where the differ the most -> conlusion to how to improve models? \\
	

	
	%%%%%%%%%%%%%%%%%%%%%%%%%%%%%%%%%%%%%%%%%%%%%%%%%%%%%%%%%%%%%%%%%%%%%%%%%%%%%%%%%%%%%%%%%%%%%%%%%%%%
	\newpage
	\pagestyle{empty}
	\printbibliography
	
\end{document}

